\begin{proof}[Proof of \cref{obj:lem:published-prognostic-capacity}]
\begin{enumerate}
\item Fix \(d\geq2\) and \(0<s\leq1\). Every law in \(\mathcal Q_{d,s}\) has prognosis equal almost everywhere to a member of \(\mathcal H_{d,s}\), so transitivity of almost-everywhere equality gives the forward inclusion of the stated images.

For the reverse inclusion, use the Gaussian completion in \cref{obj:synth_4}. Its single-unit law can be written as
\[
\begin{gathered}
\tau(u)=0.2+0.1\sqrt2\sin(2\pi u_1),
\qquad u\in[0,1]^d,\\
\operatorname{Law}(U,\varepsilon_0,\varepsilon_1)
=\operatorname{Unif}([0,1]^d)\otimes
\mathcal N(0,1)\otimes\mathcal N(0,1),\\
P_m=\operatorname{Law}\!\left(
U,\,
m(U)-\frac{\tau(U)}2+\varepsilon_0,\,
m(U)+\frac{\tau(U)}2+\varepsilon_1
\right),
\qquad m\in\mathcal H_{d,s}.
\end{gathered}
\]
The function \(m\) is square integrable, \(\tau\) is bounded, and both Gaussian errors are square integrable. Closure of \(L^2\) under finite linear combinations therefore gives
\[
\mathbb E_{P_m}[O(0)^2]<\infty,
\qquad
\mathbb E_{P_m}[O(1)^2]<\infty.
\]
The displayed construction is a conditional outcome law given \(U\). Its conditional midpoint mean is
\[
m_{P_m}(u)
=m(u)+\frac{\mathbb E[\varepsilon_0]+\mathbb E[\varepsilon_1]}2
=m(u)
\quad\text{for almost every }u.
\]
Consequently \(P_m\in\mathcal Q_{d,s}\). Every null representative of \(m\) is thus a null representative of \(m_{P_m}\), proving the reverse image inclusion and the asserted saturation.
% lean: CausalSmith.Experimentation.BivariateSobolevDesignCapacity.publishedPrognosticImage_eq_sobolevImage

\item Define the Gaussian subclass and the two-point witness by
\[
\begin{gathered}
\mathcal Q^{\mathrm G}_{d,s}
=\{P_m:m\in\mathcal H_{d,s}\},\\
\operatorname{Law}(U,\zeta_{\mathrm{tp}})
=\operatorname{Unif}([0,1]^d)\otimes
\frac{\delta_{-1}+\delta_1}{2},
\qquad
P_{\mathrm{tp}}=\operatorname{Law}(U,\zeta_{\mathrm{tp}},\zeta_{\mathrm{tp}}).
\end{gathered}
\]
The preceding step gives
\[
\mathcal Q^{\mathrm G}_{d,s}\subseteq\mathcal Q_{d,s}.
\]
For the witness,
\[
m_{P_{\mathrm{tp}}}(u)=\mathbb E[\zeta_{\mathrm{tp}}]=0,
\qquad
\mathbb E_{P_{\mathrm{tp}}}[O(0)^2]
=\mathbb E_{P_{\mathrm{tp}}}[O(1)^2]=1.
\]
The zero function has zero canonical components and zero pooled component budget, hence belongs to \(\mathcal H_{d,s}\). Thus \(P_{\mathrm{tp}}\in\mathcal Q_{d,s}\).

The outcome diagonal distinguishes this law from every Gaussian completion:
\[
P_{\mathrm{tp}}\{O(0)=O(1)\}=1.
\]
Indeed, in the Gaussian construction, conditioning on \(U\) and \(\varepsilon_0\) gives
\[
\Pr_{P_m}\!\left(O(0)=O(1)\mid U,\varepsilon_0\right)
=
\Pr\!\left(\varepsilon_1=\varepsilon_0-\tau(U)\mid U,\varepsilon_0\right)
=0,
\]
because the independent Gaussian error \(\varepsilon_1\) assigns probability zero to every singleton. Integrating yields
\[
P_m\{O(0)=O(1)\}=0.
\]
Therefore
\[
P_{\mathrm{tp}}\notin\mathcal Q^{\mathrm G}_{d,s},
\qquad
\mathcal Q^{\mathrm G}_{d,s}\subsetneq\mathcal Q_{d,s}.
\]
% lean: CausalSmith.Experimentation.BivariateSobolevDesignCapacity.gaussianLawClass_ssubset_frozenUniformClass

\item\label{proof:published-prognostic-capacity:excess} Fix \(n,d\geq2\) and \(\pi\in\mathcal D_{n,d,s}\). We first handle its almost-everywhere fairness. For a deterministic covariate array \(\mathbf x\in([0,1]^d)^n\), define
\[
\pi^{\mathrm{fair}}(\mathbf x)=
\begin{cases}
\pi(\mathbf x),
&
\displaystyle
\sum_{z\in\mathcal Z_n}z_i\pi(\mathbf x,\{z\})=0
\quad\text{for every }i,\\[4pt]
\displaystyle
\bigotimes_{i=1}^n\frac{\delta_{-1}+\delta_1}{2},
&\text{otherwise}.
\end{cases}
\]
Each displayed conditional mean is a finite Borel sum. This defines a Borel probability kernel that is fair at every covariate array and satisfies
\[
\pi^{\mathrm{fair}}(\mathbf x)=\pi(\mathbf x)
\quad\text{for }
\operatorname{Unif}([0,1]^d)^{\otimes n}\text{-almost every }\mathbf x.
\]

A covariate-based kernel also supplies conditional assignment independence. More explicitly, for an arbitrary covariate law and pre-assignment outcome kernel with domains
\[
\mu_X\in\mathcal P\bigl(([0,1]^d)^n\bigr),
\qquad
\kappa:([0,1]^d)^n\longrightarrow\mathcal P(\mathbb R^{n\times2}),
\]
the experiment is
\[
\operatorname{Law}(X,Y,Z)(d\mathbf x,dY,dz)
=\mu_X(d\mathbf x)\,\kappa(\mathbf x,dY)\,\pi(\mathbf x,dz).
\]
Its conditional distributions satisfy
\[
\operatorname{Law}(Z\mid X,Y)=\pi(X)
=\operatorname{Law}(Z\mid X),
\qquad
\operatorname{Law}(Y,Z\mid X)=\kappa(X)\otimes\pi(X).
\]
The first equality follows from the construction; integrating out \(Y\) gives the second conditional distribution. Their equality gives conditional independence, also for \(\pi^{\mathrm{fair}}\).

For any single-unit law \(P\) with bounded strictly positive covariate density and finite outcome second moments, the independent covariate sample law satisfies
\[
\operatorname{Law}_{P}\bigl((U_i)_{i=1}^n\bigr)
\ll \operatorname{Unif}([0,1]^d)^{\otimes n}.
\]
Hence the replacement preserves the entire experiment law:
\[
\operatorname{Law}_{P,\pi}\!\left(
(U_i,O_i(0),O_i(1))_{i=1}^n,Z\right)
=
\operatorname{Law}_{P,\pi^{\mathrm{fair}}}\!\left(
(U_i,O_i(0),O_i(1))_{i=1}^n,Z\right).
\]
Apply the published excess-variance identity
\citep[Proposition 2.3, equation (2.4), with Assumption 2.1 and Definition 2.2]{Cytrynbaum2026Limits}
to the pointwise fair replacement and transfer it through this equality. It gives
\[
\mathcal V_n(\pi,P)
=
\frac4n\,
\mathbb E_{P,\pi}
\left(\sum_{i=1}^n Z_i m_P(U_i)\right)^2
\geq0.
\]

Now take \(P\in\mathcal Q_{d,s}\), and choose \(m\in\mathcal H_{d,s}\) with \(m_P=m\) almost everywhere. Independence and uniformity of the unit covariates give
\[
\operatorname{Law}_{P,\pi}\bigl((U_i)_{i=1}^n,Z\bigr)
=
\operatorname{Unif}([0,1]^d)^{\otimes n}(dX)\,\pi(X,dZ).
\]
Moreover, the finite sample avoids the exceptional null set simultaneously:
\[
\Pr_{P,\pi}\!\left(
m_P(U_i)=m(U_i)\text{ for every }i
\right)=1.
\]
Square integrability of \(m\) makes the finite signed sum square integrable. Substituting the preceding two identities yields
\[
0\leq\mathcal V_n(\pi,P)
=\mathcal L_{n,d,s}(\pi,m_P)
=\mathcal L_{n,d,s}(\pi,m).
\]
This also proves invariance of the signing loss under the null representatives used here.
% lean: CausalSmith.Experimentation.BivariateSobolevDesignCapacity.published_uniform_excess_eq_loss

\item For each fixed admissible \(\pi\), the forward prognostic-image inclusion and the preceding identity give
\[
\sup_{P\in\mathcal Q_{d,s}}\mathcal V_n(\pi,P)
\leq
\sup_{m\in\mathcal H_{d,s}}\mathcal L_{n,d,s}(\pi,m).
\]
Conversely, the image equality supplies, for every \(m\in\mathcal H_{d,s}\), a law \(P\in\mathcal Q_{d,s}\) with \(m=m_P\) almost everywhere. For that law,
\[
\mathcal L_{n,d,s}(\pi,m)
=\mathcal V_n(\pi,P)
\leq
\sup_{P\in\mathcal Q_{d,s}}\mathcal V_n(\pi,P).
\]
Taking the function supremum proves
\[
\sup_{P\in\mathcal Q_{d,s}}\mathcal V_n(\pi,P)
=
\sup_{m\in\mathcal H_{d,s}}\mathcal L_{n,d,s}(\pi,m).
\]
% lean: CausalSmith.Experimentation.BivariateSobolevDesignCapacity.publishedWorst_eq_of_image_and_excess

\item The kernel construction in \cref{proof:published-prognostic-capacity:excess} establishes the independence requirement of \cref{obj:ass:assignment-independence} for every Borel probability assignment kernel. Thus a published-admissible kernel, which already satisfies fairness under the original uniform sample, satisfies both requirements of \cref{obj:def:design-class}.

Conversely, a kernel in that design class has the required fairness. In the experiment generated from independent copies of any single-unit law, drawing signs from \(\pi(X)\) gives
\[
\operatorname{Law}\!\left(
Z\mid X,(O_i(0),O_i(1))_{i=1}^n
\right)
=\pi(X)
=\operatorname{Law}(Z\mid X).
\]
This proves the published conditional-independence requirement throughout its stated law class. Consequently,
\[
\pi\text{ satisfies published admissibility}
\quad\Longleftrightarrow\quad
\pi\in\mathcal D_{n,d,s}.
\]
% lean: CausalSmith.Experimentation.BivariateSobolevDesignCapacity.published_admissible_iff_designClass

\item\label{proof:published-prognostic-capacity:capacity} Taking the infimum of the fixed-design equality over the common design class gives
\[
\inf_{\pi\in\mathcal D_{n,d,s}}
\sup_{P\in\mathcal Q_{d,s}}\mathcal V_n(\pi,P)
=
\inf_{\pi\in\mathcal D_{n,d,s}}
\sup_{m\in\mathcal H_{d,s}}\mathcal L_{n,d,s}(\pi,m)
=R_s(n,d),
\]
where the last equality is \cref{obj:def:criterion}.
% lean: CausalSmith.Experimentation.BivariateSobolevDesignCapacity.publishedCapacity_eq_of_worst

\item\label{proof:published-prognostic-capacity:comparison} Fix \(0<s\leq1\) and density bounds \(0<\lambda\leq1\leq\Lambda<\infty\). For \(P\in\mathcal Q_{d,s}\), choose \(m\in\mathcal H_{d,s}\) with \(m_P=m\) almost everywhere. Uniform density meets the comparison bounds:
\[
\frac{dP_U}{du}=1,
\qquad
\lambda\leq1\leq\Lambda.
\]
The final second-moment assertion of \cref{obj:lem:exact-reflection-budget} supplies
\[
\int m_P(u)^2\,dP_U(u)
=
\int_{[0,1]^d}m(u)^2\,du
\leq1.
\]
The permitted intercept can be chosen as zero:
\[
m_P(u)=0+m(u)\quad\text{almost everywhere}.
\]
Together with the finite outcome second moments, these statements verify every condition for membership in \(\mathcal A_{d,s}\). Therefore
\[
\mathcal Q_{d,s}\subseteq\mathcal A_{d,s}.
\]
For each admissible design, inclusion increases the law supremum. Taking infima and using \cref{proof:published-prognostic-capacity:capacity} gives
\[
R_s(n,d)
=
\inf_{\pi\in\mathcal D_{n,d,s}}
\sup_{P\in\mathcal Q_{d,s}}\mathcal V_n(\pi,P)
\leq
\inf_{\pi\in\mathcal D_{n,d,s}}
\sup_{P\in\mathcal A_{d,s}}\mathcal V_n(\pi,P).
\]
All excesses in these classes are nonnegative by the general identity in \cref{proof:published-prognostic-capacity:excess}.
% lean: CausalSmith.Experimentation.BivariateSobolevDesignCapacity.publishedCapacity_mono_class

\item With the expected absolute-supremum contraction inequality
\citep[Theorem 12(4), with Definition 2]{BartlettMendelson2002Complexities}
as its cited input, \cref{obj:thm:full-design-lower} supplies
\[
c_s b_s(n,d)\leq R_s(n,d),
\qquad
c_s=\frac{2}{(48+32\pi^2)16384^s},
\qquad
b_s(n,d)=\min\left\{1,\left(\frac{\binom d2}{n}\right)^s\right\}.
\]
The constant calibration is
\[
16384^s\leq16384,
\qquad
0<c_1=\frac{2}{(48+32\pi^2)16384}\leq c_s.
\]
Indeed, the power inequality follows from \(s\leq1\), and the denominator is positive. Combining the lower bound with \cref{proof:published-prognostic-capacity:comparison} proves
\[
c_s b_s(n,d)
\leq R_s(n,d)
\leq
\inf_{\pi\in\mathcal D_{n,d,s}}
\sup_{P\in\mathcal A_{d,s}}\mathcal V_n(\pi,P).
\]
% lean: CausalSmith.Experimentation.BivariateSobolevDesignCapacity.cLower_uniform_positive

\item Finally, let \(d_n\geq2\) for every \(n\), and suppose the displayed bounded-density capacities tend to zero. For \(n\geq2\), define
\[
\begin{gathered}
\mathfrak C_n
=
\inf_{\pi\in\mathcal D_{n,d_n,s}}
\sup_{P\in\mathcal A_{d_n,s}}\mathcal V_n(\pi,P),\\
B_n=\binom{d_n}{2},
\qquad
q_n=\frac{B_n}{n}.
\end{gathered}
\]
These quantities are used on the tail \(n\geq2\), which determines the limit. The comparison just proved gives
\[
0\leq c_s\min\{1,q_n^s\}\leq\mathfrak C_n.
\]
Since \(\mathfrak C_n\to0\) and \(c_s>0\), squeezing and division by \(c_s\) yield
\[
b_s(n,d_n)=\min\{1,q_n^s\}\longrightarrow0.
\]
Eventually this minimum is less than one, so on that tail
\[
q_n^s=b_s(n,d_n),
\qquad
q_n=b_s(n,d_n)^{1/s}\longrightarrow0,
\]
using \(s>0\).

The exact pair count and \(d_n\geq2\) imply
\[
B_n=\frac{d_n(d_n-1)}2,
\qquad
\frac{d_n^2}{4}\leq B_n\leq\frac{d_n^2}{2}.
\]
For the lower inequality, \(d_n-1\geq d_n/2\); for the upper inequality, \(d_n-1\leq d_n\). Hence
\[
0\leq\frac{d_n^2}{n}\leq4q_n\longrightarrow0.
\]
Taking square roots on the positive-\(n\) tail gives
\[
\frac{d_n}{\sqrt n}
=\sqrt{\frac{d_n^2}{n}}\longrightarrow0,
\]
which is \(d_n=o(\sqrt n)\).
% lean: CausalSmith.Experimentation.BivariateSobolevDesignCapacity.dimension_necessity_of_scale_lower
\end{enumerate}
\end{proof}
